\section{What I take from this}
\label{sec:conclusion}

For the trigger studied here, prediction is a way to adapt a margin threshold
to the estimated rate of decline.  Its value depends on whether the extra lead
changes what recovery can accomplish and how much unnecessary early action it
avoids.  In the timing model, recovery duration and deadline slack set that
boundary, and observation age consumes the available lead one for one.

The three-robot experiment shows a practical benefit of this adaptation.  Both
tuned rules succeed in all 1,620 conflict episodes, while prediction reduces
mean completion time by $2.779$\,s (paired 95\% interval
$[2.659,2.902]$\,s).  It also cuts unnecessary clear-path serializations from
1,080 to 21 out of 1,620 episodes.  Waiting until closer to the point when
recovery is needed saves time in this fixed-order crossing setup.

The two-robot experiment shows why the rate estimate matters.  The tuned fixed
rule records $0/6750$ conflict collisions and prediction records $2/6750$, both
where the conflict is revealed late and braking is slow.  A rate estimate that
lags a steepening decline can use up the lead the predictive rule was meant to
provide.

My takeaway is that the recovery requirement is the right place to start.
When recovery needs similar lead times across cases but margins decline at
different rates, adapting the threshold can reduce unnecessary waiting.
Observation delay and changing recovery durations are the next variables I
would test, using the tuned fixed threshold as the baseline.

\emph{Code and data.}  The simulators, tests, episode-level results, and figure
scripts are available at \url{\repositoryurl}.
